\textbf{Visual Attention Structure in VLMs.}
A growing body of research examines how vision-language models (VLMs) allocate visual attention in order to assess their visual grounding capabilities. Prior studies show that, although vision tokens often receive limited overall attention, increasing attention mass alone does not improve spatial reasoning performance; instead, the \emph{localization} of attention is the dominant factor~\citep{chenspatial}. In particular, overlap between attention maps and ground-truth object regions strongly correlates with answer correctness, indicating that attention localization encodes meaningful grounding signals rather than serving as a purely interpretative tool.

Further work confirms that VLM attention is frequently mislocalized, with excessive focus on non-informative regions, degrading performance on visual grounding tasks such as counting and spatial reasoning~\citep{kangsee, anwenbin2025mitigating}. Several strategies have been proposed to mitigate this issue, including attention reallocation and masking schemes. Notably, suppressing attention to distractor regions is often more effective than amplifying attention to target regions, highlighting the importance of shaping attention distributions rather than merely increasing their magnitude~\citep{sengupta2025can}.

Finally, attention allocation varies substantially across layers and heads. Rather than being uniformly distributed, grounding-relevant attention is often concentrated in a small subset of specialized heads, which exhibit more accurate localization and improved interpretability~\citep{kang2025your, kim2025interpreting}. These findings suggest that attention mechanisms constitute a promising intervention point for improving visual grounding in VLMs.

\textbf{Supervised Representation Objectives.} 
Several lines of work explore supervised objectives that align internal model representations with external visual signals. Representation alignment (REPA) has proven effective in other generative settings, for example diffusion models, where aligning hidden states with pretrained visual encoders improves both training efficiency and generation quality~\citep{yurepresentation}.

Building on this idea, recent work aligns internal representations of VLMs with those of pretrained vision encoders, leading to improved visual grounding and multimodal reasoning~\citep{yoon2025visual}. However, such approaches operate on static visual features that are independent of the textual prompt and generated tokens, limiting their ability to capture context-dependent grounding behavior.

More closely related to our work, \citet{chefer2022optimizing} propose directly supervising attention maps in unimodal vision transformers using ground-truth region annotations, demonstrating improved robustness and domain generalization. Unlike static feature representations, attention maps in VLMs are inherently prompt- and token-dependent, evolving dynamically during text generation. This property enables context-sensitive visual grounding that cannot be captured by feature-level alignment alone. We build on this insight by extending attention-level supervision to the vision-language setting, aligning token-conditioned attention maps with region-level grounding annotations during generation.

Existing approaches either apply inference-time attention heuristics, or align static visual representations that are independent of the generated text. To our knowledge, no prior work introduces a training-time attention supervision objective that explicitly aligns token-conditioned attention in VLMs with region-level grounding annotations. Our work fills this gap by supervising dynamic attention maps during text generation, enabling context-sensitive visual grounding.

% -\citep{paiss2023teaching} - in CLIP, a counting-contrastive loss is used to align image-text representations, improving counting ability and making more interpretable attention maps. Can we find similar results in VLMs?